\lecture{8}
\title{Minimum Spanning Trees}

\begin{document}

\textbf{Definition (Tree).} A tree is a connected acyclic graph.

\textbf{Definition (Spanning Tree).} A spanning tree of a graph $G = (V, E)$ is a tree $T = (V, F)$ with $F \subseteq E$.

\begin{claim}
    Every connected graph has a spanning tree.
\end{claim}

\begin{proof}
    There are many proofs (induction, BFS, \dots).  Here's an algorithmic proof:
    \begin{quote}
        \textbf{Algorithm.} Initialize $F \leftarrow E$.  While $F$ contains a cycle $C$, repeatedly pick edges from $C$ at random to remove from $F$.  Return $F$ when done.
    \end{quote}
    This terminates and yields a connected, acyclic graph (i.e. a spanning tree), done.  Here's another proof:
    \begin{quote}
        \textbf{Algorithm.} Initialize $F \leftarrow \emptyset$.  While $(V, F)$ is not connected, repeatedly draw edges between connected components of $(V, F)$.  Return $F$ when done.
    \end{quote}
    This terminates and yields a connected, acyclic graph (i.e. a spanning tree), done.
\end{proof}

\begin{problem}{Minimum Spanning Tree}
    Find the spanning tree of a weighted graph with least total edge weight.
\end{problem}

Here's a useful fact for building spanning tree intuition, from the cycle perspective.

\begin{claim}
    Say $T = (V, F)$ is a spanning tree of $G = (V, E)$, and consider any edge $e = \overline{uv} \in E \setminus F$.  Then:
    \begin{itemize}
        \item There is a \emph{unique} path $P_e \subseteq F$ connecting $u$ and $v$ in $T$.
        \item There is a \emph{fundamental cycle of $e$ with respect to $T$} defined by $P_e \cup \{e\}$.
        \item For all $e' \in P_e$, the subgraph $T' := (V, F \setminus \{e'\} \cup \{e\})$ is also a spanning tree.
    \end{itemize}
    Note that the first two bullet points don't really require $e \in E$, or care about $G$ at all, really.
\end{claim}

And here's a similar fact, from the cutting perspective.

\begin{claim}
    Say $T = (V, F)$ is a spanning tree of $G = (V, E)$, and consider any edge $e = \overline{uv} \in F$.  Then:
    \begin{itemize}
        \item If we remove $e$ from $T$, we get two connected components $A$ and $B$ (containing $u$ and $v$, respectively).
        \item For any $e' \in E$ connecting $A$ and $B$, the subgraph $T' := T \setminus \{ e \} \cup \{e'\}$ is a spanning tree.
    \end{itemize}
\end{claim}

% With these two observations, we're ready to state our algorithm.

% \begin{solution}{Algorithm Outline}
%     Initialize $F \leftarrow \emptyset$.  While $F$ is not connected, randomly pick \emph{safe} edges $e \in E \setminus F$ to add to $F$ (for some definition of \emph{safe}).  Return $F$.
% \end{solution}

% It is necessary (and sufficient!) to have $F$ be invariably a subset of some MST throughout this algorithm.

% \begin{claim}
%     Let $F \subseteq E$ be a strict subset of some MST, so $V = S \sqcup T$ such that no edge in $F$ connects $S$ and $T$.
    
%     If $e$ is any minimal-weight edge connecting $S$ and $T$, then $F \cup \{e\}$ must be a subset of some MST.
% \end{claim}

% \begin{proof}
%     Let $T$ be an MST containing $F$; say it has some edge $e'$ connecting $S$ and $T$.

%     If $e' = e$, we're fine.  And if $e' \neq e$, then $T \setminus \{e'\} \cup \{e\}$ is also a not-heavier MST, so we're also fine.
    
%     \dots wait, why is $T \setminus \{e'\} \cup \{e\}$ still connected?  The argument goes that if $e = \overline{uv}$, then $T \cup \{e\}$ has a cycle  ~ $\blacksquare$
% \end{proof}

This inspires Kruskal's Algorithm, which constructs $T$ by adding edges to a forest.

\begin{solution}{Kruskal's Algorithm}
    Here's the algorithm, in full.
    \begin{enumerate}
        \item Initialize $F \leftarrow \emptyset$.  Initialize a $\textsc{Union-Find}$ data structure to store connected components.  Sort $E$ by weight.
        \item Scan through all edges $e_i = (u_i, v_i)$ for all $i = 1, 2, \dots, |E|$, in increasing order of weight.
        \begin{itemize}
            \item For each $(u_i, v_i)$ such that $\textsc{Find-Set}(u_i) \neq \textsc{Find-Set}(v_i)$, add $e_i$ to $F$ and apply $\textsc{Union}(u_i, v_i)$.
        \end{itemize}
        \item Finally, return $F$.
    \end{enumerate}
\end{solution}

\begin{proof}
    It suffices to demonstrate by induction that $F$ is always a subset of some MST.

    Suppose $F \subseteq T$ initially, for some MST $T$.  Now we draw an edge $e = (u_i, v_i)$ between the connected component $A$ of $u_i$ in $(V, F)$ and $B := V \setminus A$.  (Intuitively, we should think of edge $e$ as bridging across the ``cut'' $V = A \sqcup B$.)

    If $e = (u_i, v_i) \in T$, we're fine.  Otherwise, take the fundamental cycle of $e$ with respect to $T$; this cycle must bridge across the ``cut'' $V = A \sqcup B$ a second time along some edge $e'$.

    You can now check that $T' := T \setminus \{e'\} \cup \{e\}$ is also an MST, so $F \cup \{e\}$ remains a subset of an MST $T'$. ~ $\blacksquare$
\end{proof}

Here's our runtime analysis of Kruskal's Algorithm.

\begin{solution}{Kruskal's Algorithm Runtime}
There's only three major steps.
    \begin{enumerate}
        \item Initializing the $\textsc{Union-Find}$ data structure takes $O(|V|)$ time.
        \item Sorting $E$ takes $O(|E| \log |E|) = O(|E| \log |V|)$ time (since $|E| = O(|V|^2)$).
        \item All of the $\textsc{Union-Find}$ operations take $O(|E| \cdot \alpha(|V|))$ time.
    \end{enumerate}
So in total, we can say the runtime of Kruskal's Algorithm is $O(|E| \log |V|)$.
\end{solution}

This also inspires Prim's Algorithm, which grows $F$ by connecting vertices to a tree.

\begin{remark}
    Recall the following properties of a \emph{priority queue} implemented by a \emph{Fibonacci heap}.
    \begin{itemize}
        \item A priority queue $Q$ stores keys together with their values.
        \item The method $Q.\textsc{DeleteMin}()$ returns the lowest-value key of $Q$ and deletes it in $O(\log n)$ amortized time.
        \item The methods $Q.\textsc{Insert}(k, v)$ and $Q.\textsc{DecreaseKey}(k, v)$ run in $O(1)$ amortized time.
    \end{itemize}
\end{remark}

\begin{solution}{Prim's Algorithm}
    Here's the algorithm, in full:
    \begin{enumerate}
        \item Initialize $F = (V_F, E_F) \leftarrow (\emptyset, \emptyset)$.
        \item Initialize a $\textsc{Priority Queue}$ $Q$ and an array $\textsc{Parent}[u]$.
        \begin{itemize}
            \item The queue $Q$ will store keys $v \in V$ with values $d(v) := \min \{ w(u, v) \mid u \in V_F \text{ and } \overline{uv} \in E\}$.
            \item The $\textsc{Parent}$ array will store for each $v$ the vertex with $w(v, \textsc{Parent}[v]) = d(v)$.
            \item Initialize $\textsc{Parent}[u] \leftarrow \textsc{Null}$ for all $u \in V$.
        \end{itemize}
        \item Pick some starting vertex $s$, and initialize $Q.\textsc{Insert}(s, 0)$ and $Q.\textsc{Insert}(v, \infty)$ for all $v \neq s$.
        \item While $Q$ is nonempty\dots
        \begin{enumerate}
            \item Add the edge $(u, \textsc{Parent}[u])$ to $F$, where $u \leftarrow Q.\textsc{DeleteMin}()$.
            
            (Add nothing when $\textsc{Parent}[u] = \textsc{Null}$, as on the first iteration.)
            \item Now that we've added $u$ to $F$, we need to make some updates.  For each $v$ neighboring $u$ still in $Q$\dots
            \begin{itemize}
                \item Update $d(v) \leftarrow  \min\{d(v), w(u, v)\}$, and update $Q.\textsc{DecreaseKey}(v, d(v))$ as needed.
                \item If $d(v)$ was updated, then also update $\textsc{Parent}[v] \leftarrow u$.
            \end{itemize}
        \end{enumerate}
        \item Return $F$ when done.
    \end{enumerate}
\end{solution}

\begin{proof}
    Again, we wish to show $F$ is invariably a subset of some MST.

    It's actually the same argument as before.  Consider the cut $V = V_F \sqcup (V \setminus V_F)$.  You can do \emph{whatever} to bridge this cut and the result will still be an ST, minimal if you do the ``\emph{whatever}'' with minimal weight. ~ $\blacksquare$
\end{proof}

This runs in $O(|E| + |V|\log |V|)$ time, since\dots
\begin{itemize}
    \item We have $|V|$ iterations of $Q.\textsc{DeleteMin}()$, each of which takes $O(\log |V|)$ time.
    \item The updates $d(v) \leftarrow \min \{d(v), w(u, v)\}$ have amortized $O(1)$ cost with a Fibonacci heap, $O(|E|)$ total.
\end{itemize}

\end{document}